\documentclass[../../book]{subfiles}

\begin{document}

In previous schemes, the verifier relied on algebraic structures: pairings in KZG, and group homomorphisms in Pedersen commitments. Here, we rely on nothing but hash functions. We need to verify that a committed function is a low-degree polynomial by examining only a few of its evaluations. Verifying this directly is impossible: the function could perfectly match a polynomial everywhere except at a single, unchecked point.

This is precisely why we need to reformulate the problem: if we extend a polynomial over a domain much larger than its degree, any function that is not a valid polynomial will end up being 'far away' from actual polynomials. To leverage this property and solve our problem, we will turn to coding theory~\cite{guruswami23}.
%--------------------------------
\subsection{Linear Codes}
%--------------------------------
\begin{intuition}
  Suppose the sender wants to send a message $m$ over a noisy channel. The sender first uses an encoding function to map the $k$ message symbols into $n$ symbols, called a \emph{codeword}, and then sends it over the channel. The receiver gets a \emph{received word} of $n$ symbols and tries to decode it, recovering the original $k$ message symbols.
\end{intuition}

\begin{definition}[Code]\label{def:code}
  A \emph{code} of block length $n$ over an alphabet $\Sigma$ is a subset $C \subseteq \Sigma^n$, whose elements are called \emph{codewords}. We write $M := |C|$ for the number of codewords and $k := \log_{|\Sigma|} M$ for the number of message symbols. An \emph{encoding} for $C$ is an injective map $E : \Sigma^k \to \Sigma^n$ whose image is $C$.
\end{definition}

Let us define a notion of distance that captures how close two vectors $\mathbf{u}$ and $\mathbf{v}$ are.
\begin{definition}[Hamming distance]\label{def:hamming-distance}
  Given $\mathbf{u}, \mathbf{v} \in \Sigma^n$, the \emph{Hamming distance} between $\mathbf{u}$ and $\mathbf{v}$, denoted as $\Delta(\mathbf{u}, \mathbf{v})$, is defined to be the number of positions in which they differ.
\end{definition}

To more clearly understand the Hamming distance, suppose we have two vectors $\mathbf{u}=00000$ and $\mathbf{v}=10100$. Then, the Hamming distance between $\mathbf{u}$ and $\mathbf{v}$ is $\Delta(\mathbf{u}, \mathbf{v})=2$ since they differ in the 1st and 3rd positions.

Hamming distance alone does not give you the full picture. For example, two vectors of length $1000$ that differ in $5$ places are much more similar than two vectors of length $10$ that also differ in $5$ places. Because of this, it is standard practice to normalize the distance by the block length.

\begin{definition}[Relative Hamming distance]\label{def:rel-distance}
  Let $\mathbf{u}, \mathbf{v} \in \Sigma^n$. The \emph{relative Hamming distance} between $\mathbf{u}$ and $\mathbf{v}$, denoted as $\delta(\mathbf{u}, \mathbf{v})$, is defined to be the Hamming distance between $\mathbf{u}$ and $\mathbf{v}$ divided by $n$:
  \begin{equation*}
    \delta(\mathbf{u}, \mathbf{v}) = \frac{1}{n} \Delta(\mathbf{u}, \mathbf{v}) \in [0,1]
  \end{equation*}
\end{definition}

While comparing individual vectors is helpful, a code's effectiveness depends on the entire set. Specifically, we care about the distance between codewords, where the worst-case spacing between any two codewords is the defining metric.
\begin{definition}[Minimum distance]\label{def:min-distance}
  Let $C\subseteq \Sigma^n$ be a code. The \emph{minimum distance} of $C$, denoted by $d_{\min}(C)$, is defined as 
  \begin{equation*}
    d_{\min}(C) = \min_{\mathbf{c}_1\neq \mathbf{c}_2 \in C} \Delta(\mathbf{c}_1, \mathbf{c}_2)
  \end{equation*}
    and the \emph{relative minimum distance} of $C$ is $\mu(C) := \frac{1}{n} d_{\min}(C)$.
\end{definition}

The minimum distance is the foundation of error correction. If $d_{\min}(C) = d$, introducing fewer than $d$ errors will never result in another valid codeword. Furthermore, with fewer than $d/2$ errors, the corrupted string remains strictly closer to the original codeword than to any other. Essentially, a larger minimum distance forces codewords further apart, so the original codeword can be recovered as long as fewer than $d/2$ errors occur.

While calculating $d_{\min}(C)$ directly from the definition requires comparing every possible pair, practical codes offer a much faster shortcut. To introduce it, we first need to define the Hamming weight.

\begin{definition}[Hamming weight]\label{def:hamming-weight}
  The \emph{Hamming weight} of a vector $\mathbf{c} \in \Sigma^n$, denoted by $w(\mathbf{c})$, is defined as the number of non-zero symbols in $\mathbf{c}$.
\end{definition}

When $\Sigma$ is a field, these two concepts are directly linked by the equation $\Delta(\mathbf{u}, \mathbf{v}) = w(\mathbf{u} - \mathbf{v})$, which means that two vectors differ precisely in the coordinates where their difference is non-zero. 

In practice, cryptographic codes are never just arbitrary subsets of $\Sigma^n$ --- they are built with underlying algebraic structure. Assuming a code is closed under addition and scalar multiplication imposes virtually no practical limitations, yet it provides a twofold advantage.

\begin{definition}[Linear code]\label{def:linear-code}
  Let $\mathbb{F}_q$ be a finite field. A code $C \subseteq \mathbb{F}_q^n$ is a \emph{linear code} if it is a linear subspace of $\mathbb{F}_q^n$. If $C$ has dimension $k$ and minimum distance $d$, we denote it as an $[n,k,d]_q$ code.
\end{definition}

The first payoff is the shortcut promised above. Since the difference of two codewords of a linear code is again a codeword, the minimisation over all pairs collapses into a minimisation over single codewords.

\begin{proposition}\label{prop:dist-eq-weight}
  Let $C \subseteq \mathbb{F}_q^n$ be an $[n,k,d]_q$ code. Then
  \begin{equation*}
    d = \min_{\mathbf{0} \neq \mathbf{c} \in C} w(\mathbf{c}).
  \end{equation*}
\end{proposition}

\begin{proof}
  First, we show that $d$ is at most the minimum weight. Let $\mathbf{c}'$ be a non-zero codeword of minimum weight. Since $C$ is a linear subspace, $\mathbf{0} \in C$, and the distance from $\mathbf{c}'$ to $\mathbf{0}$ equals its weight. Hence $d \leq \Delta(\mathbf{0}, \mathbf{c}') = w(\mathbf{c}')$.

  Next, we show that $d$ is at least the minimum weight. Let $\mathbf{c}_1 \neq \mathbf{c}_2 \in C$ be such that $\Delta(\mathbf{c}_1, \mathbf{c}_2) = d$. By linearity, $\mathbf{c}_1 - \mathbf{c}_2 \in C$, and it is non-zero since $\mathbf{c}_1 \neq \mathbf{c}_2$. The non-zero symbols of $\mathbf{c}_1 - \mathbf{c}_2$ occur exactly in the positions where the two codewords differ, so $w(\mathbf{c}_1 - \mathbf{c}_2) = \Delta(\mathbf{c}_1, \mathbf{c}_2) = d$. Therefore the minimum weight of a non-zero codeword is at most $d$.
\end{proof}

The second payoff is that the dimension $k$ of a linear code is exactly the number of field elements each codeword carries, since $|C| = q^k$. This lets us put a price on the redundancy.

\begin{definition}[Rate]\label{def:rate}
  Let $C\subseteq \Sigma^n$ be a code. The \emph{rate} of $C$ is defined as 
  \begin{equation*}
    \rho = \frac{k}{n}\in [0,1]
  \end{equation*}
\end{definition}

A codeword consists of $n$ symbols but carries only $k$ symbols of information; thus $\rho$ is the fraction of each codeword that is actual payload. The rate and the relative minimum distance pull in opposite directions: adding redundancy lowers the rate but makes room for codewords that lie further apart.

\begin{proposition}[Singleton bound]\label{prop:singleton}
  Every $[n,k,d]_q$ code satisfies $k \leq n-d+1$, equivalently $|C| \leq q^{n-d+1}$.
\end{proposition}

Codes attaining this bound are optimal and belong to a special class known as maximum distance separable.

\begin{definition}[MDS code]\label{def:mds}
  An $[n,k,d]_q$ code is called \textbf{maximum distance separable} (MDS) if $d = n-k+1$,
  that is, if it attains the Singleton bound.
\end{definition}
%--------------------------------
\subsection{Reed-Solomon Code}
%--------------------------------
We now introduce the primary code that forms the foundation for the rest of this section. Notably, it attains the Singleton bound.
\begin{definition}[Reed-Solomon Code]\label{def:rs}
  Let $\mathbb{F}_q$ be a finite field. Let $\bm{\alpha} = (\alpha_0, \alpha_1, \ldots, \alpha_{n-1})$ be a sequence of $n$ distinct elements (also called evaluation points) of $\mathbb{F}_q$ and choose $n$ and $k$ such that $k\leq n \leq q$. \emph{The Reed-Solomon} code is an encoding function $\mathsf{RS}_q[\bm{\alpha},k]: \mathbb{F}_q^k \to \mathbb{F}_q^n$ defined as follows: a message $\mathbf{m} = (m_0, m_1, \ldots, m_{k-1}) \in \mathbb{F}_q^k$ is mapped to a polynomial $\mathbf{m}\rightarrow f_{\mathbf{m}}(T)$, where
  \begin{equation*}
    f_{\mathbf{m}}(T) = \sum_{i=0}^{k-1} m_i\cdot T^i
  \end{equation*}
  Note that $f_{\mathbf{m}}(T)$ is a polynomial of degree at most $k-1$. The encoding of $\mathbf{m}$ is the evaluation of $f_{\mathbf{m}}(T)$ at all the $\alpha_i$'s:
\begin{equation*}
    \mathsf{RS}_q[\bm{\alpha},k](\mathbf{m}) = \big(f_{\mathbf{m}}(\alpha_0), f_{\mathbf{m}}(\alpha_1), \ldots, f_{\mathbf{m}}(\alpha_{n-1})\big).
\end{equation*}
\end{definition}

When $\bm{\alpha}$ and $k$ are clear from context, we simply write $\mathsf{RS}$.

\Cref{tab:rs-example} lists all codewords of the $[3,2,2]_3$ Reed--Solomon code whose evaluation points are $\mathbb{F}_3 = \{0,1,2\}$. The second row shows the polynomial $f_{\mathbf{m}}(T)$ corresponding to each message $\mathbf{m} \in \mathbb{F}_3^2$, and the codewords are ordered by the coefficients of these polynomials.

\begin{table}[h!]
\centering
\scriptsize
\begin{tabular}{|l|c|c|c|c|c|c|c|c|c|}
\hline
$\mathsf{RS}(\mathbf{m})$ & (0,0,0) & (1,1,1) & (2,2,2) & (0,1,2) & (1,2,0) & (2,0,1) & (0,2,1) & (1,0,2) & (2,1,0) \\ \hline
$f_{\mathbf{m}}(T)$       & $0$ & $1$ & $2$ & $T$ & $T+1$ & $T+2$ & $2T$ & $2T+1$ & $2T+2$ \\ \hline
\end{tabular}
\caption{All codewords of the $[3,2,2]_3$ Reed-Solomon code with evaluation points $\mathbb{F}_3$, together with the corresponding message polynomials.}
\label{tab:rs-example}
\end{table}

\begin{proposition}\label{prop:rs-linear}
  The Reed-Solomon code $\mathsf{RS}$ is a linear code.
\end{proposition}

\begin{proof}
    To prove this, let us consider $a\in \mathbb{F}_q$ and $f(T), g(T)\in \mathbb{F}_q[T]$ be two polynomials of degree at most $k-1$. Then, we have that $af(T)$ and $f(T)+g(T)$ are also polynomials of degree at most $k-1$. Thus, let messages $\mathbf{m}_1, \mathbf{m}_2 \in \mathbb{F}_q^k$ be mapped to $f_{\mathbf{m}_1}(T), f_{\mathbf{m}_2}(T) \in \mathbb{F}_q[T]$ are polynomials of degree at most $k-1$. Then, it is easy to verify that:

    Let $a \in \mathbb{F}_q$ and $\mathbf{m}_1, \mathbf{m}_2 \in \mathbb{F}_q^k$. The polynomials $f_{\mathbf{m}_1}(T)$ and $f_{\mathbf{m}_2}(T)$ have degree at most $k-1$, and so do $a f_{\mathbf{m}_1}(T)$ and $f_{\mathbf{m}_1}(T) + f_{\mathbf{m}_2}(T)$. Since coefficients are added coordinate-wise:
    \begin{equation*}
      f_{\mathbf{m}_1}(T) + f_{\mathbf{m}_2}(T) = f_{\mathbf{m}_1+\mathbf{m}_2}(T),
    \end{equation*}
    and 
    \begin{equation*}
      a f_{\mathbf{m}_1}(T) = f_{a\mathbf{m}_1}(T)
    \end{equation*}

    In other words, the property of linearity holds for Reed-Solomon codes:
    \begin{gather*}
      \mathsf{RS}(\mathbf{m}_1) + \mathsf{RS}(\mathbf{m}_2) = \mathsf{RS}(\mathbf{m}_1+\mathbf{m}_2), \\ a \mathsf{RS}(\mathbf{m}_1) = \mathsf{RS}(a\mathbf{m}_1)
    \end{gather*}
    Therefore $\mathsf{RS}$ is an $[n,k]_q$ linear code.
\end{proof}

\begin{proposition}\label{prop:rs-distance}
  The Reed-Solomon code $\mathsf{RS}$ is an $[n,k,n-k+1]_q$ code. That is, it attains the Singleton bound.
\end{proposition}
\begin{proof}
  Let us fix arbitrary distinct $\mathbf{m}_1, \mathbf{m}_2\in \mathbb{F}_q^k$. Note that $f_{\mathbf{m}_1}(T), f_{\mathbf{m}_2}(T)\in \mathbb{F}_q[T]$ are distinct polynomials of degree at most $k-1$ since $\mathbf{m}_1\neq \mathbf{m}_2$. Then, the polynomial $f_{\mathbf{m}_1}(T)-f_{\mathbf{m}_2}(T)$ is also a non-zero polynomial of degree at most $k-1$. 

  Note that $w(\mathsf{RS}(\mathbf{m}_1)-\mathsf{RS}(\mathbf{m}_2))=\Delta(\mathsf{RS}(\mathbf{m}_1), \mathsf{RS}(\mathbf{m}_2))$. The weight of $\mathsf{RS}(\mathbf{m}_1)-\mathsf{RS}(\mathbf{m}_2)$ is $n$ minus the number of zeroes in $\mathsf{RS}(\mathbf{m}_1)-\mathsf{RS}(\mathbf{m}_2)$, which is equal to $n$ minus the number of roots that $f_{\mathbf{m}_1}(T)-f_{\mathbf{m}_2}(T)$ has in $\{\alpha_0, \alpha_1, \ldots, \alpha_{n-1}\}$. That is,
  \begin{equation*}
    \Delta(\mathsf{RS}(\mathbf{m}_1), \mathsf{RS}(\mathbf{m}_2)) = n - \#\{i \in [n] : f_{\mathbf{m}_1}(\alpha_i)=f_{\mathbf{m}_2}(\alpha_i)\}
  \end{equation*}
  Since $f_{\mathbf{m}_1}(T)-f_{\mathbf{m}_2}(T)$ is a non-zero polynomial of degree at most $k-1$, by \Cref{cor:root-count} it has at most $k-1$ roots in $\mathbb{F}_q$.
  Therefore, we have that the weight of $\mathsf{RS}(\mathbf{m}_1)-\mathsf{RS}(\mathbf{m}_2)$ is at least $n-(k-1)=n-k+1$. Therefore $d\geq n-k+1$, and since \Cref{prop:singleton} implies that $d\leq n-k+1$, we have that $d=n-k+1$.

  The argument above also shows that the distinct polynomials $f_{\mathbf{m}_1}(T)$ and $f_{\mathbf{m}_2}(T)$ are mapped to distinct codewords $\mathsf{RS}(\mathbf{m}_1)$ and $\mathsf{RS}(\mathbf{m}_2)$. Therefore, the code contains $q^k$ codewords, and has dimension $k$.
\end{proof}

\begin{remark}
  By \Cref{prop:rs-distance}, Reed-Solomon codes are MDS codes. In particular,
  $\mu(\mathsf{RS}) = 1 - \rho + \frac{1}{n}$, so lowering the rate directly increases the relative distance.
\end{remark}

%--------------------------------
\subsection{Proximity to a Code}\label{subsec:proximity}
%--------------------------------

The concepts discussed so far belong to classical coding theory, where the main objective is to recover the original message from a corrupted codeword. In our context, however, the requirement is weaker: the verifier does not actually need to know \emph{which} polynomial the prover committed to, but only whether the committed vector is close to \emph{some} valid codeword. To formalize this, we evaluate the distance between an arbitrary vector and the code as an entire set.

\begin{definition}[Distance to a code]\label{def:dist-to-code}
    Let $C \subseteq \Sigma^n$ be a code and $\mathbf{u} \in \Sigma^n$. The distance from
    $\mathbf{u}$ to $C$ is
    \begin{equation*}
        \Delta(\mathbf{u}, C) := \min_{\mathbf{c} \in C} \Delta(\mathbf{u}, \mathbf{c}),
    \end{equation*}
    and the relative distance is $\delta(\mathbf{u}, C) := \frac{1}{n}\Delta(\mathbf{u}, C)$.
\end{definition}

\begin{definition}[$\delta$-far]\label{def:delta-far}
    For $\delta \in [0,1]$, a vector $\mathbf{u} \in \Sigma^n$ is \textbf{$\delta$-far} from $C$
    if $\delta(\mathbf{u}, C) > \delta$, and \textbf{$\delta$-close} otherwise.
\end{definition}

Why is this notion useful? Assume $\mathbf{u}$ is $\delta$-far from $C$. This implies that regardless of which codeword $\mathbf{c}$ we compare it to, the two will differ in strictly more than $\delta n$ coordinates. Consequently, a single uniformly random query will detect a discrepancy with a probability greater than $\delta$. If we make $t$ independent queries, the probability of missing the error entirely drops below $(1-\delta)^t$. Thus, a constant relative distance ensures that the soundness error decays exponentially with the number of queries. The fundamental challenge, however, is that the verifier does not know in advance which codeword $\mathbf{c}$ they are checking against. Constructing an efficient test that overcomes this limitation using only a small number of queries to $\mathbf{u}$ is exactly what the FRI protocol achieves.

%--------------------------------
\subsection{Reed-Solomon Proximity Problem}\label{subsec:rs-proximity}
%--------------------------------

We now state the problem that FRI solves~\cite{bensasson18}. Let $D \subseteq \mathbb{F}_q$ be a set of $n$ evaluation points, and let $\mathsf{RS}_q[D,k]$ denote the Reed-Solomon code whose evaluation points are the elements of $D$ (\Cref{def:rs}). The verifier has \emph{oracle access} to a function $f : D \to \mathbb{F}_q$: it cannot read $f$ as a whole, but it may query its value at any point of $D$.
In practice, the prover commits to $f$ with a Merkle tree (\Cref{section:commitment-schemes}), so that each query costs one opening.

\begin{definition}[Reed-Solomon proximity problem]\label{def:rs-proximity}
  Given oracle access to $f : D \to \mathbb{F}_q$ and $\delta \in (0,1)$, distinguish, with high confidence and few queries, between the case $f \in \mathsf{RS}_q[D,k]$ and the case that $f$ is $\delta$-far from $\mathsf{RS}_q[D,k]$.
\end{definition}

When the verifier is restricted to oracle access to $f$, the problem is known as \emph{Reed-Solomon proximity testing}. In this model, exactly $k+1$ queries are both necessary and sufficient: the verifier reads $k$ points to interpolate the unique polynomial of degree less than $k$, and then checks it for consistency against one additional random point. However, since $k$ is typically extremely large in real-world applications, this straightforward interpolation is computationally infeasible.

To overcome this limitation, we allow the prover to participate. In \emph{Reed-Solomon proximity verification}, the prover and verifier interact via an IOP (\Cref{fig:iop}). During each round, the prover commits to an additional oracle, and the verifier is allowed to query any of them. This interactive setup is known as an \emph{IOP of proximity} (IOPP), making standard proximity testing just the degenerate case where the prover provides no auxiliary information.

The \emph{Fast Reed-Solomon IOPP} (FRI) is precisely this type of protocol. It achieves linear prover complexity in $n$, while keeping both the verifier's computational effort and the query complexity strictly logarithmic in $n$.

%--------------------------------
\subsection{FRI protocol}\label{subsec:fri}
%--------------------------------
Unlike the polynomial commitment schemes from \Cref{section:commitment-schemes}, FRI is not formalized as a static tuple of algorithms (such as $\mathsf{Setup}$ and $\mathsf{Com}$). Instead, it operates as a two-phase interactive protocol, as outlined in \Cref{fig:fri-scheme}.

\begin{figure}[h!]
    \centering
    \scalebox{0.8}{
    \begin{tikzpicture}[
        arr/.style   = {-{Stealth[length=2.5mm]}, line width=0.4mm},
        act/.style   = {draw=oc-blue-7, fill=oc-blue-1, rounded corners=2pt,
                        line width=0.4mm, inner sep=4pt, font=\small},
        data/.style  = {font=\small, inner sep=2pt},
        phase/.style = {draw=gray, dashed, rounded corners=8pt, line width=0.4mm},
        halo/.style  = {preaction={draw, white, line width=2.2mm}},
    ]
        % --- Verifier column
        \draw[fill=gray!20!white, ultra thick, rounded corners=4pt] (8.6,-3.9) rectangle (10.4,6.0);
        \node[inner sep=0pt, align=center] at (9.5,7.0)
            {\includegraphics[width=1.0cm]{contents/images/common/verifier.png}\\ Verifier $\mathcal{V}$};

        % --- Prover
        \node[inner sep=0pt, align=center] at (0,7.0)
            {\includegraphics[width=1.0cm]{contents/images/common/prover.png}\\ Prover $\mathcal{P}$};

        % --- Phase frames
        \draw[phase] (-1.9,0.55) rectangle (8.3,6.0);
        \draw[phase] (-1.9,-3.9) rectangle (8.3,0.15);
        \node[font=\scriptsize\scshape, text=gray] at (6.9,5.6) {Commit phase};
        \node[font=\scriptsize\scshape, text=gray] at (6.9,-0.25) {Query phase};

        % --- Commit phase
        \node[data] (poly)  at (0,5.45) {polynomial $f$};
        \node[act]  (eval)  at (0,4.55) {evaluate on $D_0$};
        \node[data] (cw)    at (0,3.65) {codeword $f_i$ on $D_i$};
        \node[act]  (merk)  at (3.3,2.85) {Merkle commit};
        \node[data] (root)  at (6.3,2.85) {Merkle root};
        \node[act]  (fold)  at (0,2.0) {split-and-fold};
        \node[data] (beta)  at (6.3,2.0) {challenge $\beta_i$};
        \node[data] (final) at (0,1.1) {final value $f_r$};

        \draw[arr] (poly) -- (eval);
        \draw[arr] (eval) -- (cw);
        \draw[arr] (cw.south east) -- (merk.north west);
        \draw[arr] (merk) -- (root);
        \draw[arr] (root) -- (8.6,2.85);
        \draw[arr] (8.6,2.0) -- (beta);
        \draw[arr] (beta) -- (fold);
        \draw[arr] (fold) -- (cw);
        \draw[arr] (fold) -- (final);
        \draw[arr] (final) -- (8.6,1.1);

        % --- Query phase
        \node[data] (trees) at (3.3,-0.7) {Merkle trees};
        \node[act]  (open)  at (3.3,-1.75) {open};
        \node[data] (idx)   at (6.3,-1.75) {query points $s$};
        \node[data] (vals)  at (1.4,-2.8) {values};
        \node[data] (paths) at (5.0,-2.8) {authentication paths};

        \draw[arr, halo] (merk) -- (trees);
        \draw[arr] (trees) -- (open);
        \draw[arr] (8.6,-1.75) -- (idx);
        \draw[arr] (idx) -- (open);
        \draw[arr] (open) -- (vals);
        \draw[arr] (open) -- (paths);
        \draw[arr, rounded corners=4pt] (vals.south) |- (8.6,-3.45);
        \draw[line width=0.4mm, rounded corners=4pt] (paths.south) |- (5.6,-3.45);
    \end{tikzpicture}
    }
    \caption{Overview of the FRI protocol. In the commit phase, the prover repeatedly commits to the
    current codeword and folds it using the verifier's challenge, until a final value remains.
    In the query phase, the verifier asks to open the committed codewords at random points and checks
    the authentication paths. Adapted from~\cite{szepieniec2021anatomy}.}
    \label{fig:fri-scheme}
\end{figure}

To see how the protocol works, we first describe its main building block.

\subsubsection{Split-and-Fold scheme}
FRI follows a split-and-fold strategy: a claim is split into two claims of half the size, the verifier sends a random challenge, and the prover merges the two claims into one using this challenge. Then, after logarithmically many steps, the claim has been reduced to one of trivial size. If the original claim is false, then with high probability the final claim is false as well.

In the FRI case, this claim asserts that the given codeword corresponds to a polynomial of low degree. Let $n$ be the length of the codeword, and let $f(T)=\sum_{i=0}^{k-1} c_i\cdot T^i$ be a polynomial of degree less than $k$.

Following the divide-and-conquer strategy of the fast Fourier transform, we split a polynomial $f$ into its even and odd parts:
\begin{gather*}
  f(T) = f_E(T^2) + T \cdot f_O(T^2),
  \qquad
  f_E(T^2) = \frac{f(T) + f(-T)}{2},
  \\
  f_O(T^2) = \frac{f(T) - f(-T)}{2T}.
\end{gather*}
Indeed, the odd terms cancel in $f(T) + f(-T)$ and the even terms cancel in $f(T) - f(-T)$.
If $\deg f < k$ for an even $k$, then both $f_E$ and $f_O$ have degree less than $k/2$.
Given a random challenge $\beta \in \mathbb{F}_q$ from the verifier, the prover derives a codeword for
\begin{equation*}
  f^\star(T) := f_E(T) + \beta \cdot f_O(T)
\end{equation*}
from the codeword for $f(T)$.

Let $D$ be a subgroup of even order $n$ of the multiplicative group of the field, and let $\omega$ be a generator of this subgroup. Let $\{f(\omega^i)\}_{i=0}^{n-1}$ be the codeword for $f(T)$, corresponding to evaluation on $D$. Let $D^\star=\langle \omega^2\rangle$ be another domain of half the length, and let $\{f_E(\omega^{2i})\}_{i=0}^{n/2-1}$, $\{f_O(\omega^{2i})\}_{i=0}^{n/2-1}$ and $\{f^\star(\omega^{2i})\}_{i=0}^{n/2-1}$ be the codewords for $f_E(T)$, $f_O(T)$ and $f^\star(T)$, respectively, corresponding to evaluation on $D^\star$. As a result, we can rewrite the definition of $f^\star(T)$:
\begin{gather*}
  \{f^\star(\omega^{2i})\}_{i=0}^{n/2-1} = \left\{f_E(\omega^{2i})+\beta\cdot f_O(\omega^{2i})\right\}_{i=0}^{n/2-1} = \\
  \left\{\frac{f(\omega^i)+f(-\omega^i)}{2}+\beta \cdot\frac{f(\omega^i)-f(-\omega^i)}{2\omega^i}\right\}_{i=0}^{n/2-1} = \\
  \left\{2^{-1}\cdot \big((1+\beta\cdot \omega^{-i})\cdot f(\omega^i)+(1-\beta\cdot \omega^{-i})\cdot f(-\omega^i)\big)\right\}_{i=0}^{n/2-1}.
\end{gather*}

Note that, since $\omega^{n/2}=-1$, we have $f(-\omega^i)=f(\omega^{n/2+i})$. This makes it clear that even though the index iterates over half the range, all $n$ points of the codeword for $f(T)$ are involved.

At this point it is possible to describe one round of FRI: the prover commits to $f(T)$ by sending the Merkle root of its codeword to the verifier. The verifier responds with the random challenge $\beta$. The prover computes $f^\star(T)$ and commits to it by sending the Merkle root of $\{f^\star(\omega^{2i})\}_{i=0}^{n/2-1}$ to the verifier.

In turn, having received the two commitments, the verifier can check that they are correctly related, and rejects the proof if for some $x \in D$
\begin{equation}\label{eq:fri-fold}
  f^\star(x^2)\neq 2^{-1}\cdot \big((1+\beta x^{-1})\cdot f(x)+(1-\beta x^{-1})\cdot f(-x)\big).
\end{equation}
To do this, the verifier randomly samples an index $i \leftarrowS [n/2]$, sets $x := \omega^i$, and thus defines $3$ points:
\begin{itemize}
  \item $A:\quad (\omega^i, f(\omega^i))$,
  \item $B:\quad (\omega^{n/2+i}, f(\omega^{n/2+i}))$,
  \item $C:\quad (\beta, f^\star(\omega^{2i}))$.
\end{itemize}
Upon receiving the index $i$ from the verifier, the prover provides the $y$-coordinates along with their Merkle authentication paths. The verifier checks these paths against the corresponding roots and then verifies that $A$, $B$ and $C$ lie on a straight line. This is called the \emph{collinearity check}.

%Add exercise A and B on the same line.

After the first round is complete, the verifier and the prover return to the starting point. The prover wishes to establish that a given Merkle root decommits to a codeword whose defining polynomial has a bounded degree. However, as a result of the previous round, the length of the codeword and the number of possible non-zero coefficients of the polynomial have both been halved. Now they set $f := f^\star$ and $D := D^\star$ and repeat the process with a fresh challenge. If $k = 2^r$, then after $r = \log_2 k$ rounds of FRI, the prover and the verifier end up with a constant polynomial whose codeword is also constant, so the prover simply sends this constant to the verifier.

\begin{definition}[FRI protocol]\label{def:fri}
  Let $f_0 := f$ be given on $D_0 := D = \langle \omega \rangle$ of order $n$, with $\deg f < k = 2^r$.
  \begin{itemize}
    \item \textbf{Commit phase.} For $i = 0, 1, \ldots, r-1$:
    \begin{algoen}
      \item the prover sends the Merkle root of the codeword of $f_i$ on $D_i$;
      \item the verifier sends a random challenge $\beta_i \leftarrowS \mathbb{F}_q$;
      \item the prover computes $f_{i+1} := f_{i,E} + \beta_i \cdot f_{i,O}$ on $D_{i+1} := \{x^2 : x \in D_i\}$.
    \end{algoen}
    Finally, the prover sends the constant $f_r$.
    \item \textbf{Query phase.} Repeat $t$ times:
    \begin{algoen}
      \item the verifier samples $x_0 \leftarrowS D_0$ and sets $x_{i+1} := x_i^2$ for $i \in [r]$;
      \item for every $i \in [r]$, the prover opens $f_i(x_i)$ and $f_i(-x_i)$ together with their
      authentication paths;
      \item the verifier checks all paths and \eqref{eq:fri-fold} for every $i \in [r]$, where
      $f_r(x_r)$ is the constant sent by the prover.
    \end{algoen}
  \end{itemize}
  The verifier accepts if and only if all checks pass.
\end{definition}

\begin{remark}
  As described, FRI is interactive. In STARKs it is made non-interactive using the Fiat-Shamir
  transformation (\Cref{subsection:fiat-shamir}): each challenge $\beta_i$ and each starting point $x_0$
  is derived by hashing the Merkle roots sent so far.
\end{remark}

\subsection{Arithmetic Intermediate Representation (AIR)}\label{subsec:air}
Similar to SNARK, the STARK proof system involves intermediate stages, the first two of which are the \emph{arithmetic constraint system} and the \emph{IOP}.

\begin{definition}\label{def:air}
  The \emph{arithmetic intermediate representation (AIR)} or arithmetic internal representation, is a way of describing a computation in terms of an execution trace that satisfies a number of constraints induced by the correct evolution of the state.
\end{definition}

\begin{definition}\label{def:aet}
  Let $\mathbb{F}_q$ be the field of definition, and the computation describes the evolution of a state of $w$ registers for $\tau$ cycles. Then we define the \emph{algebraic execution trace (AET)} is the table of $\tau \times w$ field elements where every row describes the state of the system at the given point in time, and every column tracks the value of the given register.
\end{definition}

\begin{definition}\label{def:st-function}
  A \emph{state transition function} $\phi:\mathbb{F}_q^{w}\rightarrow\mathbb{F}_q^{w}$ determines the state at the next cycle as a function of the state at the previous cycle. Furthermore, a set of boundary conditions 
  \begin{equation*}
    \mathcal{B}\subseteq \mathbb{Z}_\tau \times\mathbb{Z}_w\times\mathbb{F}_q
  \end{equation*}
  fixes the values of some registers at some cycles: a triple $(i,w,e) \in \mathcal{B}$ means that register $w$ holds the value $e$ at cycle $i$.
\end{definition}

The \emph{computational integrity claim} consists of the state transition function and the boundary conditions. The \emph{witness} to this claim is the algebraic execution trace. The claim is \emph{true} if there is a witness $W\in \mathbb{F}_q^{\tau \times w}$ such that:
\begin{itemize}
  \item For every cycle, the state evolves correctly: $\forall i\in [\tau-1], f(W_{[i, :]})=W_{[i+1, :]}$
  \item All boundary conditions are satisfied: $\forall (i, w, e)\in \mathcal{B}, W_{[i, w]}=e$
\end{itemize}
Actually, the state transition function hides a lot of complexity and requires to be describable as low degree polynomials that are independent of the cycle.
However, this list of polynomials does not need to compute the next state from the current one; it merely needs to distinguish correct evolutions from incorrect ones. Specifically, the state transition function $\phi$ is represented by a list of polynomials $\mathbf{p}(x_0,\ldots,x_{w-1}, y_0,\ldots,y_{w-1})$ such that $\phi(\mathbf{x})=\mathbf{y}$ if and only if $\mathbf{p}(\mathbf{x},\mathbf{y})=0$. Say there are $m$ such state transition verification polynomials. Then the transition constraint become:
\begin{equation*}
  \forall i\in[\tau-1],\forall j\in[m]: p_j(W_{[i, 0]},\ldots,W_{[i, w-1]},W_{[i+1, 0]},\ldots,W_{[i+1, w-1]})=0
\end{equation*}
This representation admits non-determinism: the verification polynomials may have much lower degree than any polynomial that computes the next state. For instance, the transition $y = x^{-1}$ requires computing $x^{q-2}$, but it is verified by the degree-$2$ polynomial $xy - 1$.

\subsubsection{Trace interpolation}
We saw that the arithmetic constraint system described above represents the computational integrity claim as a bunch of polynomials, where each such polynomial corresponds to a constraint. However, we need to extend terms of polynomials to the witness for transformation this constraint system into an IOP, and extending the notion of valid witnesses to witness polynomials.

Let $H = \langle h \rangle$ be a multiplicative subgroup of $\mathbb{F}_q^{\times}$ of order $\tau$, which we call the \emph{trace evaluation domain}. For simplicity, we assume that $\tau$ is a power of two; otherwise the trace is padded. The trace domain is smaller than the FRI evaluation domain $D$ by a factor exactly equal to the blowup factor $1/\rho$.

For every register $j \in [w]$, the \emph{trace polynomial} $t_j \in \mathbb{F}_q[T]$ is the unique polynomial of degree less than $\tau$ such that
\begin{equation*}
  t_j(h^i) = W_{[i,j]} \quad \forall i \in [\tau].
\end{equation*}
In practice, it is computed with the inverse NTT (\Cref{section:ntt}). The trace polynomials
$t_0, \ldots, t_{w-1}$ represent the algebraic execution trace in terms of univariate polynomials.

Translating the computational integrity claim to the trace polynomials, we get:
\begin{itemize}
  \item All boundary constraints are satisfied: $\forall (i,w,e)\in \mathcal{B}, t_w(h^i)=e$.
  \item For all cycles, all transition constraints are satisfied: $\forall i\in [\tau-1],\forall j\in[m]$
  \begin{equation*}
    p_j(t_0(h^i),\dots,t_{w-1}(h^i),t_0(h^{i+1}),\dots,t_{w-1}(h^{i+1}))=0
  \end{equation*}
\end{itemize}
If we look at the left-hand side of the equation, we can see that it corresponds to a univariate polynomial $p_j(t_0(T),\dots,t_{w-1}(T),t_0(h\cdot T),\dots,t_{w-1}(h\cdot T))$. Which just means that all $m$ of these transition polynomials evaluate to $0$ in $H$.

Then we can build the following high-level protocol:
This observation gives rise to the following high-level protocol:
\begin{algoen}
  \item The prover commits to the trace polynomials $t_0(T), \ldots, t_{w-1}(T)$.
  \item The verifier checks that $t_j(h^i) = e$ for all $(i,j,e) \in \mathcal{B}$.
  \item The prover commits to the transition polynomials
  \begin{equation*}
    c_\ell(T) := p_\ell\big(t_0(T), \ldots, t_{w-1}(T), t_0(h \cdot T), \ldots, t_{w-1}(h \cdot T)\big),
    \quad \ell \in [m].
  \end{equation*}
  \item The verifier checks that $c_\ell(T)$ and $t_j(T)$ are correctly related by:
  \begin{itemize}
    \item choosing a random point $z \leftarrowS \mathbb{F}_q^{\times}$;
    \item querying the values $t_j(z)$ and $t_j(h \cdot z)$ $\forall j \in [w]$;
    \item evaluating the transition verification polynomials $p_\ell$ at
    \begin{equation*}
      \big(t_0(z), \ldots, t_{w-1}(z), t_0(h \cdot z), \ldots, t_{w-1}(h \cdot z)\big)
    \end{equation*}
    \item querying the values $c_\ell(z)$;
    \item checking that the values obtained in the previous two steps match.
  \end{itemize}
  \item The verifier checks that the transition polynomials $c_\ell(T)$ evaluate to zero on
  $\{h^i : i \in [\tau-1]\}$.
\end{algoen}

In fact, the commitment to the transition polynomials can be omitted. Instead, the verifier uses
the values of $t_j$ at $z$ and $h \cdot z$ to compute the value of $c_\ell(T)$ at the one point
needed to verify that $c_\ell(T)$ evaluates to $0$ on $\{h^i : i \in [\tau-1]\}$.



\end{document}